\chapter*{Sammendrag}

Dette prosjektet utvikler en numerisk metode for å identifisere flere samtidige lekkasjer i et rør, der lekkasjeidentifikasjon omfatter både lokalisering og estimering av lekkasjestørrelse. Metoden bygger utelukkende på stasjonære målinger av trykkhøyde og vannføring ved innløp og utløp og er ment å støtte programvarebasert lekkasjeidentifikasjon i vannforsyningssystemer med begrenset måleutstyr. Lekkasjer i rørledninger fører til betydelige vanntap, økte driftskostnader og økt risiko for infrastrukturen, noe som gjør praktiske metoder for lekkasjeidentifikasjon viktige. Eksisterende stasjonære metoder er imidlertid i stor grad begrenset til én eller to lekkasjer, mens metoder for flere lekkasjer vanligvis krever transiente data eller ekstra sensorer.\\

For å håndtere dette gapet ble lekkasjeidentifikasjonsproblemet formulert som et ulineært residualproblem og løst ved hjelp av en Levenberg--Marquardt-metode med en analytisk utledet Jakobimatrise. Metoden ble evaluert gjennom et omfattende simulert testrammeverk og mot publiserte simulerte eksperimenter med to lekkasjer. Resultatene viser at metoden fungerer pålitelig for to og tre lekkasjer og at den i mange tilfeller fortsatt kan gi nyttig lokalisering for fire lekkasjer. En kondisjonstallsanalyse viste videre at det inverse problemet raskt blir dårligere kondisjonert når antall lekkasjer øker. Sammen med de empiriske resultatene tyder dette på at identifisering av fem eller flere samtidige lekkasjer trolig ligger nær den øvre grensen for hva som er numerisk gjennomførbart ved bruk av kun stasjonære endepunktsmålinger.\\

Funnene viser at stasjonære endepunktsmålinger kan identifisere flere samtidige lekkasjer enn sammenlignbare publiserte stasjonære metoder. Samtidig avdekker de en praktisk grense som skyldes den iboende kondisjoneringen til det inverse problemet.